\item Un calorímetro de aluminio, con una masa de $100 \text{ g}$, contiene $250 \text{ g}$ de agua. El calorímetro y el agua están en equilibrio térmico a $10.0 ^\circ\text{C}$. Dos bloques metálicos se colocan en el agua. Uno es un trozo de cobre de $50.0 \text{ g}$ a $80.0 ^\circ\text{C}$. El otro tiene una masa de $70.0 \text{ g}$ y originalmente está a una temperatura de $100 ^\circ\text{C}$. Todo el sistema se estabiliza a una temperatura final de $20.0 ^\circ\text{C}$. (a) Determine el calor específico de la muestra desconocida. (b) Con los datos de la tabla 7.1, ¿puede hacer una identificación positiva del material desconocido? ¿Puede identificar un material posible? (c) Explique sus respuestas al inciso (b).
